\SourceLeaf{402}%
\par
If \(t\) is the complex number for the centre of gravity, then \(3t=\alpha+\beta+\gamma\) and hence
\[
\begin{aligned}
3t-2u&=\alpha+(\beta-\gamma)\operatorname{cotg}A.i\\
&=\beta+(\gamma-\alpha)\operatorname{cotg}B.i\\
&=\gamma+(\alpha-\beta)\operatorname{cotg}C.i
\end{aligned}
\]
\par
This \(3t-2u\) is the complex number for the point where the three perpendiculars from the vertices to the opposite sides meet.

Subtracting therefore gives
\[
\begin{aligned}
0&=\alpha-\beta+\{\alpha\operatorname{cotg}B+\beta\operatorname{cotg}A-\gamma(\operatorname{cotg}A+\operatorname{cotg}B)\}.i\\
0&=\beta-\gamma+\{\beta\operatorname{cotg}C+\gamma\operatorname{cotg}B-\alpha(\operatorname{cotg}B+\operatorname{cotg}C)\}.i\\
0&=\gamma-\alpha+\{\gamma\operatorname{cotg}A+\alpha\operatorname{cotg}C-\beta(\operatorname{cotg}C+\operatorname{cotg}A)\}.i
\end{aligned}
\]
\par\medskip\Needspace{5\baselineskip}
\pdfbookmark[2]{Second handwritten remark.}{carnot-handschriftlich-ii}
\begin{center}[\textit{Second handwritten remark.}]\end{center}
\hspace*{\parindent}Let \(a,b,c,d\) be four points on the circumference of a circle of radius \(1\) and also the complex numbers corresponding to those points [with the complex number corresponding to the centre taken to be \(0\)]. Let \(p,q,r\) be the intersections of the lines \(\tfrac{ab}{cd},\tfrac{ac}{bd},\tfrac{ad}{bc}\). Finally, let \(p^*\) be the midpoint of the chord at whose endpoints the tangents meet at \(p\), and let \(q^*,r^*\) be defined in the same way. Then, with primed letters always denoting the respective conjugates, we have
\[
\begin{aligned}
p&=\frac{abc+abd-acd-bcd}{ab-cd}=a-\frac{b(a-c)(a-d)}{ab-cd}\\
q&=\frac{abc+acd-abd-bcd}{ac-bd}=a-\frac{c(a-b)(a-d)}{ac-bd}\\
r&=\frac{acd+abd-abc-bcd}{ad-bc}=a-\frac{d(a-b)(a-c)}{ad-bc}
\end{aligned}
\]
\[
p-q=(a-d)(b-c)\frac{abd+acd-abc-bcd}{(ab-cd)(ac-bd)}=\frac{(a-d)(b-c)(ad-bc)r}{(ab-cd)(ac-bd}
\]
\[
p^*=\frac1{p'}=\frac{ab-cd}{a+b-c-d}=a-\frac{(a-c)(a-d)}{a+b-c-d}
\]
\begingroup\footnotesize
\[
p^*-q=\frac{(a-d)(b-c)(ad-bc)}{(ac-bd)(a+b-c-d)}=p^*.\frac{(a-d)(b-c)(ad-bc)}{(ab-cd)(ac-bd)}=p^*\frac{p-q}{r}
\]
\endgroup
or \(p^*(q+r-p)=qr\), likewise \(q^*(p+r-q)=pr\), and \(r^*(p+q-r)=pq\)
